\subsubsection{Base presentations and multiplier ideals}

The following constructions take place in the category of generalized
data specified above. In particular, the whole adjoint is a
Bott--Chern class, whereas the boundary and nef traces on a singular
space are interpreted as currents. On a smooth carrier the structure
boundary is a real divisor; its negative components may be exceptional
over the space of the pair. The nef part is specified on that carrier
and pulled back to higher models.

\begin{lemma}[Adjoint-preserving preparation]
\label{prog:base-presentation-preparation}
Let $(X,B+\mathbf M)$ be a generalized pair on a normal compact
K\"ahler space, with globally modified-big total boundary
$B+\mathbf M_X$. Write
\[
 A=[K_X+B+\mathbf M_X],\qquad
 \mathcal J_o=\mathcal J(X,B+\mathbf M).
\]
There are generalized data on $X$ with the same adjoint class, a
projective carrier $p:U\to X$ with $U$ smooth and compact K\"ahler, an SNC real
divisor $C$ whose negative components are $p$-exceptional, and a
K\"ahler class $\beta$ such that
\begin{equation}\label{prog:base-presentation-prepared}
 p^*A=[K_U+C]+\beta,\qquad
 \mathcal J_1=p_*\mathcal O_U(-\lfloor C\rfloor)
                 \supseteq\mathcal J_o.
\end{equation}
The carrier may be required to dominate any prescribed projective
modification of $X$, with the structure and exceptional divisors
having simultaneous SNC support. The prepared total boundary is
globally modified big. In particular, the preparation preserves the
gklt condition.
\end{lemma}

\begin{proof}
Choose a common projective log resolution $p_0:U_0\to X$ of the
original nef carrier and the prescribed modification, and write
\[
 p_0^*A=[K_{U_0}+C^0]+\mathbf M_{U_0}.
\]
The class $\mathbf M_{U_0}$ is nef, and the negative part $(C^0)^-$
is exceptional over $X$. The modified-big preparation
\cite[Lemmas~2.10 and~2.23]{HX} can be made with the prescribed
resolution. We give its divisor and ideal details. Choose an effective divisor $E_0$ with sufficiently large coefficients
and full $p_0$-exceptional support so that the class
\[
 [(C^0)^+]+\mathbf M_{U_0}+[E_0]
\]
is big. Resolve a Kodaira decomposition of this class. On a further
projective modification $a:U\to U_0$ with $U$ smooth one obtains
\[
 a^*\bigl([(C^0)^+]+\mathbf M_{U_0}+[E_0]\bigr)
       =\kappa+[D],\qquad D\geq0,
\]
where $\kappa$ is K\"ahler. The support can be made SNC together
with all the other divisors. On any additional blowup, subtract a
sufficiently small effective exceptional class from the pulled-back
K\"ahler class and add it to the divisor; an exceptional divisor
whose negative is relatively ample makes the resulting class
K\"ahler again.

Put $p=p_0a$ and
\[
 C^{\mathrm{alt}}=D-a^*(C^0)^--a^*E_0-K_{U/U_0}.
\]
Then
\[
 p^*A=[K_U+C^{\mathrm{alt}}]+\kappa,
\]
and every negative component of $C^{\mathrm{alt}}$ is
$p$-exceptional. Let $C^{\mathrm{old}}$ be the crepant structure
boundary of the original data on this same model. For $0<t\ll1$,
set
\[
 C=(1-t)C^{\mathrm{old}}+tC^{\mathrm{alt}},\qquad
 \beta=(1-t)a^*\mathbf M_{U_0}+t\kappa.
\]
The class $\beta$ is K\"ahler, the adjoint identity is unchanged,
and negative coefficients still occur only on exceptional divisors.
There are only finitely many coefficients. Thus $t$ can be chosen so
that
\[
 \lfloor C\rfloor\leq\lfloor C^{\mathrm{old}}\rfloor.
\]
Pushing forward the corresponding inclusion of invertible sheaves
gives~\eqref{prog:base-presentation-prepared}. The K\"ahler class
$\beta$ plus the strict transform of the effective trace boundary is
big; its pushforward is the prepared total boundary. If the original pair
is gklt, its multiplier ideal is the unit ideal, and hence so is
$\mathcal J_1$.
\end{proof}

\begin{proposition}[Base presentations for projective contractions]
\label{prog:base-presentation}
Let $f:X\to Z$ be a projective contraction between normal irreducible
compact K\"ahler spaces. Suppose that $(X,B+\mathbf M)$ has globally
modified-big total boundary and that
\[
 A=[K_X+B+\mathbf M_X]=f^*\gamma,
 \qquad \gamma\in H^{1,1}_{\mathrm{BC}}(Z).
\]
There is an adjoint-preserving preparation as in
Lemma~\ref{prog:base-presentation-preparation}, with multiplier ideal
$\mathcal J_1\supseteq\mathcal J_o$, for which
\begin{equation}\label{prog:base-presentation-Stein}
 R^if_*\mathcal J_1=0\quad(i>0),\qquad
 f_*\mathcal O_{N_1}=\mathcal O_Z/(f_*\mathcal J_1),
 \qquad N_1=V(\mathcal J_1).
\end{equation}
If $N_1$ dominates $Z$, then $f_*\mathcal J_1=0$.
Otherwise there are generalized data $(Z,B_Z+\mathbf N)$, a
projective carrier $\tau:T\to Z$ with $T$ smooth and compact K\"ahler, a real
SNC structure divisor $C_T$, and a K\"ahler class $\beta_T$ such that
\begin{equation}\label{prog:base-presentation-identities}
 \begin{gathered}
 [K_Z+B_Z+\mathbf N_Z]=\gamma,\qquad
 \tau^*\gamma=[K_T+C_T]+\beta_T,\qquad
 \mathbf N_T=\beta_T,\\
 B_Z=\tau_*C_T\geq0,\qquad
 B_Z+\mathbf N_Z\ \text{is globally modified big}.
 \end{gathered}
\end{equation}
The negative components of $C_T$ are $\tau$-exceptional, and
\begin{equation}\label{prog:base-presentation-ideal}
 \mathcal J_Z:=\mathcal J(Z,B_Z+\mathbf N)
       =\tau_*\mathcal O_T(-\lfloor C_T\rfloor)
       \supseteq f_*\mathcal J_1.
\end{equation}
In particular the base data are gklt wherever $f_*\mathcal J_1$ is
the unit ideal. If the source data are gklt, the base data are gklt.
All adjoint identities are preserved by pullback to higher models.
\end{proposition}

\begin{proof}
Choose a projective resolution $\mu:W\to Z$, with $W$ smooth
compact K\"ahler and reduced exceptional divisor $F$ having SNC
support. Apply Lemma~\ref{prog:base-presentation-preparation} on a
carrier dominating the main component of the base change. We obtain
projective maps $p:U\to X$ and $q:U\to W$ with $fp=\mu q$ and
\begin{equation}\label{prog:base-presentation-source}
 [K_U+C]+\beta=p^*A=q^*\mu^*\gamma,
 \qquad \mathcal J_1=p_*\mathcal O_U(-\lfloor C\rfloor).
\end{equation}
Here $U$ is smooth compact K\"ahler, $\beta$ is K\"ahler, and all
negative components of $C$ are $p$-exceptional. The maps are
projective by composition and factorization over the separated base
$W$. The map $q$ has connected fibres: this holds over the open set
where $\mu$ is an isomorphism, and its Stein factorization is finite
and bimeromorphic over the normal space $W$.

The integral divisor $-\lfloor C\rfloor$ has class
\[
 [-\lfloor C\rfloor]
   =[K_U+\{C\}]+\beta-(fp)^*\gamma.
\]
The pair $(U,\{C\})$ is klt with SNC boundary. Relative
Kawamata--Viehweg vanishing, first for $p$ and then for $fp$, applies
because the positive class modulo pullbacks is K\"ahler
\cite[Theorem~2.2]{HX}. Equivalently, the real line bundle
$-K_U-C$ is relatively ample,
so the vanishing follows from \cite[Theorem~1.2]{FujinoVanishing2025}
with the identity as its first morphism. Leray gives $R^if_*\mathcal J_1=0$ for
$i>0$. Applying $f_*$ to the sequence defining $N_1$, and using
$f_*\mathcal O_X=\mathcal O_Z$, proves
\eqref{prog:base-presentation-Stein}. If $N_1$ dominates $Z$, the
ideal $f_*\mathcal J_1\subseteq\mathcal O_Z$ vanishes generically
and therefore vanishes. This proves that alternative.

Suppose from now on that $N_1$ does not dominate $Z$. Every
$q$-horizontal coefficient of $C$ is then less than one. Indeed,
a component with coefficient at least one forces the multiplier
ideal to be proper along its image in $X$, and its image would make
$N_1$ dominate $Z$. If $Z$ is a point, the asserted base presentation
is the dimension-zero presentation. We therefore assume
$b=\dim Z>0$.

\paragraph{The fibre metric.}
Choose a smooth representative $\rho$ of
$\mu^*\gamma-c_1(K_W)$, a K\"ahler form $\omega_W$, and a
K\"ahler form representing $\beta$, still denoted by $\beta$.
For some $a>0$,
\[
 \beta\geq a q^*\omega_W.
\]
Consider the holomorphic line bundle
\[
 L=\mathcal O_U(-K_{U/W}-\lfloor C\rfloor).
\]
Equation~\eqref{prog:base-presentation-source} gives a singular
Hermitian metric $h$ on $L$ whose curvature current is
\begin{equation}\label{prog:base-presentation-curvature}
 \Theta_h(L)=\beta+[\{C\}]-q^*\rho.
\end{equation}
Its weights are smooth functions plus the divisorial logarithms
with coefficients in $\{C\}$. We normalize $dd^c$ so that
$dd^c\log|z|^2=[z=0]$ and write squared-norm weights as $e^{-\varphi}$.
Let $e$ be the canonical meromorphic section of
\[
 K_{U/W}\otimes L=\mathcal O_U(-\lfloor C\rfloor),
\]
so that $\operatorname{div}(e)=-\lfloor C\rfloor$.

Choose an analytic Zariski open set $W^\circ\subseteq W$ on which
$q$ is a submersion, the horizontal boundary is relatively SNC,
there are no vertical boundary components, $\mu$ is an isomorphism
onto an open subset of $Z$, and coherent base change holds in degree
zero. On $q^{-1}(W^\circ)$, the divisor $-\lfloor C\rfloor$ is
effective and exceptional over $X$. Normality and the contraction
property therefore identify
\[
 q_*(K_{U/W}\otimes L)|_{W^\circ}
       =\mathcal O_{W^\circ}e.
\]
Here $K_{U/W}$ is viewed as the relative canonical line bundle. The
fibre norm
\[
 I(w)=\int_{U_w}|e_w|_h^2
\]
is positive and finite. The fractional SNC weights are integrable;
any zeros of $e_w$ only improve integrability. Relative SNC
coordinates and a partition of unity show that $I$ is smooth:
parameter derivatives act on smooth factors multiplying fixed
integrable monomial weights.

Set $\psi=-\log I$ on $W^\circ$. Then
\begin{equation}\label{prog:base-presentation-open-positivity}
 \rho+dd^c\psi\geq a\omega_W.
\end{equation}
Here are the details for the singular weights. Locally choose
$\ell$ with $dd^c\ell=\rho-a\omega_W$, and put
$h'=he^{-q^*\ell}$. Its weights are psh, and its fibre norm satisfies
$I'=e^{-\ell}I$. Restrict the family to a small coordinate line disk
$\Delta_R$, with centre $0$. The sharp $L^2$ extension theorem gives
a twisted top form extending $e_0$ and satisfying
\begin{equation}\label{prog:base-presentation-sharp-extension}
 \int_{q^{-1}(\Delta_R)}|E|_{h'}^2
        \leq\pi R^2 I'(0).
\end{equation}
Top-form norms here are normalized to use Euclidean area in the
base coordinate. We use the Stein hypersurface form of the optimal
extension theorem \cite[Theorem~2.2 and \S3.5, (3.1)]{GuanZhouOptimal2015},
with defining function the disk coordinate divided by $R$.

To check applicability, embed the restricted projective family into
projective space over the disk and remove a relative hyperplane
section that does not contain the central fibre. The complement is
Stein. A holomorphic section of $L$ on this Stein space can be chosen
nonzero at a point of the central fibre; deleting its zero divisor
trivializes $L$ and leaves a Stein space. The central fibre is smooth
and connected, hence irreducible, so both deletions leave a dense
open subset of it. For a zero-dimensional fibre choose the deletions
to avoid its point. In the resulting trivialization the weight is
psh, and the sharp extension theorem applies to canonical forms.
The extension extends across the deleted hypersurfaces in reverse
order by local $L^2$ removability: in every original frame the metric
is bounded below by a smooth positive metric. The extension retains
its norm bound and its prescribed restriction to the central fibre.

After dividing $E$ by the base differential, the rank-one direct
image identifies it with $b(t)e$, where $b$ is holomorphic and
$b(0)=1$. Thus the disk average of $|b(t)|^2I'(t)$ is at most $I'(0)$.
Jensen's inequality and the submean inequality for $\log|b|^2$ give
\[
 \log I'(0)
 \geq\operatorname{avg}_{\Delta_R}
             \bigl(\log|b(t)|^2+\log I'(t)\bigr)
 \geq\operatorname{avg}_{\Delta_R}\log I'(t).
\]
The same argument applies to every sufficiently small disk at every
centre and in every complex line direction. Hence $-\log I'=\psi+\ell$ is psh, which proves
\eqref{prog:base-presentation-open-positivity}.

\paragraph{Extension over the base.}
Let $D$ be a prime divisor in $W\setminus W^\circ$. Choose a component
$Q$ of $q^*D$ dominating $D$. When $D$ is not $\mu$-exceptional, choose
$Q$ not $p$-exceptional: take a component over the corresponding prime
on $Z$ already on $X$ and then its strict transform. Put
\[
 u=\operatorname{ord}_Q(q^*D),\qquad c=\operatorname{coeff}_Q(C).
\]
At a general point of $Q$, avoiding all other relevant divisors,
there are local coordinates with
\[
 t_1=z_1^u,\qquad t_j=z_j\quad(2\leq j\leq b).
\]
Indeed the map $Q\to D$ is submersive at a general smooth point,
and a local unit in the first equation is absorbed by a holomorphic
$u$th root. Such points can be chosen above every point of a dense
analytic open subset of $D$ by generic smoothness and properness.
Integration over a fixed small polydisk in the remaining fibre
coordinates gives
\begin{equation}\label{prog:base-presentation-divisor-bound}
 I(t)\geq c_*|t_1|^{-2(c+u-1)/u}
\end{equation}
for some $c_*>0$ near each such base point. The section and metric
contribute the exponent $c$, while division by the base differential
in the relative canonical form contributes $u-1$.

For a nonexceptional $D$ the chosen $Q$ has $c\geq0$, so
\eqref{prog:base-presentation-divisor-bound} bounds $\psi$ from above
near its general points. For the finitely many exceptional divisors,
the same estimate supplies a nonnegative integer $m$ such that, if
$k$ is a smooth metric on $\mathcal O_W(F)$ and $s_F$ its divisor
section,
\[
 \varphi=\psi+m\log|s_F|_k^2
\]
is locally bounded above near general points of every divisor.
Adding local potentials of $\rho+m c_1(k)-a\omega_W$ to
$\varphi$ gives psh weights by
\eqref{prog:base-presentation-open-positivity}. These weights
extend across the divisors and then across the remaining analytic
set of codimension at least two, giving
\begin{equation}\label{prog:base-presentation-positive-current}
 S=\rho+m c_1(k)+dd^c\varphi\geq a\omega_W
\end{equation}
on all of $W$.

\paragraph{A K\"ahler carrier.}
Regularize the current in
\eqref{prog:base-presentation-positive-current} in its class, retaining
a lower bound for the potentials. More precisely, choose an
approximation with analytic singularities such that
\begin{equation}\label{prog:base-presentation-regularization}
 S'=\rho+m c_1(k)+dd^c\varphi',\qquad
 S'\geq\frac a2\omega_W,\qquad
 \varphi'\geq\varphi-O(1).
\end{equation}
These simultaneous properties are the Bergman regularization
properties of \cite{DemaillyRegularization1992}; see also
\cite[pp.~59--60]{PopoviciRegularization2004}.
Resolve the analytic singularities and the support of $F$ by a
sequence of blowups with smooth centres $r:T\to W$, with simultaneous
SNC support, and write
\[
 r^*S'=[H]+\zeta,\qquad H\geq0.
\]
The residual current $\zeta$ has locally bounded potentials. Removing
the divisorial part from the positive current
$r^*S'-(a/2)r^*\omega_W$ shows that
\[
 \zeta\geq\frac a2r^*\omega_W.
\]
The difference has vanishing Lelong numbers, so its class is nef by
regularization. Choose an effective $r$-exceptional divisor $G$ with
$-G$ relatively ample. For $0<\eta\ll1$, the class
$(a/2)r^*\omega_W-\eta[G]$ is K\"ahler. Therefore
\begin{equation}\label{prog:base-presentation-carrier-data}
 \begin{split}
 \tau&=\mu r,\qquad
 \beta_T=[\zeta]-\eta[G],\\
 C_T&=H+\eta G-mr^*F-K_{T/W}
 \end{split}
\end{equation}
has K\"ahler $\beta_T$. If $r$ is an isomorphism, omit $G$ and its
correction. The choice of $\eta$ will be decreased once more below.

The exact class calculation is
\[
 \begin{split}
 [K_T+C_T]+\beta_T
 &=r^*c_1(K_W)+[H]+[\zeta]-m[r^*F]\\
 &=r^*(c_1(K_W)+[\rho])=\tau^*\gamma.
 \end{split}
\]
All possibly negative terms in $C_T$ are $\tau$-exceptional.
Consequently $B_Z=\tau_*C_T$ is effective. Specify the nef b-data by
$\mathbf N_T=\beta_T$; the displayed identity defines generalized
data on $Z$ with adjoint class $\gamma$. At the level of currents,
pushing the identity gives the required trace identity on $Z$:
exceptional divisor currents push to zero, and a $dd^c$ difference
pushes to a $dd^c$ difference. Canonical traces can be computed using
local meromorphic canonical frames; a change of frame changes the
corresponding Chern potentials by the usual logarithmic transition
term. Thus only the whole adjoint, rather
than its individual traces, needs to be a Bott--Chern class. The
class $\beta_T$ plus the strict transform of $B_Z$ is big on $T$,
so its pushforward $B_Z+\mathbf N_Z$ is globally modified big.
These are~\eqref{prog:base-presentation-identities}.

\paragraph{The multiplier ideal.}
Take a local holomorphic function $v\in f_*\mathcal J_1$. By the
pushforward description of $\mathcal J_1$,
\begin{equation}\label{prog:base-presentation-source-orders}
 \operatorname{div}(q^*\mu^*v)\geq\lfloor C\rfloor.
\end{equation}
Thus $q^*\mu^*v\,e$, tensored with a local base canonical frame,
is a holomorphic $L$-valued top form on $U$. Its squared norm is
locally integrable by the fractional SNC weights in
\eqref{prog:base-presentation-curvature}. Over the submersion locus,
fibre integration followed by base integration is precisely this
top-form integral. Properness and a finite coordinate covering
therefore give
\begin{equation}\label{prog:base-presentation-base-integral}
 \int |\mu^*v|^2e^{-\psi}<\infty
\end{equation}
locally on $W$; the integral is taken where the original fibre
metric is defined.

The lower bound on $\varphi'$ in
\eqref{prog:base-presentation-regularization} preserves this
integrability after replacing $\psi$ by
$\varphi'-m\log|s_F|_k^2$. Change variables through $r$. Along a
prime divisor of $T$, let $h,j,k_D,d$ be respectively the coefficients
of $H$, $\operatorname{div}(\tau^*v)$, $K_{T/W}$, and $r^*F$.
The local integrand has exponent
$2(j-h+md+k_D)$, up to a bounded positive factor. Its integrability
implies $j-h+md+k_D>-1$, or, since $j$ is integral,
\[
 j\geq\lfloor h-md-k_D\rfloor.
\]
Therefore
\[
 \operatorname{div}(\tau^*v)
       \geq\lfloor H-mr^*F-K_{T/W}\rfloor.
\]
Choose $\eta>0$ sufficiently small that adding $\eta G$ does not
change any of these finitely many round-downs. The resulting
inequality is exactly
$\operatorname{div}(\tau^*v)\geq\lfloor C_T\rfloor$, proving
\eqref{prog:base-presentation-ideal} as an inclusion of ideal sheaves.
In particular, if $f_*\mathcal J_1$ is the unit ideal on an open set,
then so is $\mathcal J_Z$, and the base data are gklt there.

Finally, the nef data descend to $T$, and the structure boundaries
on higher models are defined by crepant pullback of
\eqref{prog:base-presentation-identities}. On every common refinement
of the source and base carriers, both adjoints are the pullback of
$\gamma$. Thus the identities remain exact on higher models.
\end{proof}

\begin{lemma}[Contractions of closed subspaces]
\label{prog:base-emc}
Let $D$ be a compact complex space, possibly nonreduced, and let
$\delta\in H^2(D,\R)$ have nonnegative degree on compact curves.
Use the definitions of a Moishezon morphism and of a contraction
defined by a class in \cite[Definitions~2.32 and~2.34]{HX}.
Then the following hold.
\begin{enumerate}
\item If $\delta$ is endowed with a Moishezon contraction, so is
its restriction to every closed analytic subspace of $D$.
\item Let $g:D\to V$ be a projective contraction of compact
complex spaces, so $g_*\OO_D=\OO_V$, and suppose
$\delta=g^*\epsilon$ for a class $\epsilon\in H^2(V,\R)$
having nonnegative degree on compact curves. Then $\delta$ is
endowed with a Moishezon contraction if and only if $\epsilon$
is endowed with one.
\end{enumerate}
\end{lemma}

\begin{proof}
Write $h:D\to S$ for a Moishezon contraction defined by
$\delta$. For a closed subspace $D_0\subset D$, take the Stein
factorization
\[
                 D_0\longrightarrow S_0\longrightarrow S
\]
of its restriction. The first map is a contraction, and the second
is finite. Thus a curve in $D_0$ is contracted by the first map
exactly when its $\delta$-degree is zero. Choose a projective
surjection $P\to D$ witnessing that $h$ is Moishezon, with
$P\to S$ projective. The closed subspace
$P_0=P\times_D D_0$ is projective over $S$ and projective
and surjective over $D_0$. It is also projective over $S_0$:
a relatively ample line for $P_0\to S$ remains ample on the
fibres of $P_0\to S_0$. This proves the first assertion,
including its scheme structure.

For the descent direction of the second assertion, every curve
in a fibre of $g$ has $\delta$-degree zero. A connected
projective complex space is connected by chains of compact curves
after passing to its reduction. Therefore $h$ is constant on
the underlying set of each fibre of $g$. It factors as
$h=k\circ g$ for a holomorphic map $k:V\to S$.
To see that this factorization holds also for nonreduced spaces,
choose a coordinate neighbourhood in $S$ of the image of a fibre.
Properness permits restriction to a neighbourhood of its point
in $V$ whose inverse image maps into that coordinate
neighbourhood. The coordinate functions of $h$ then descend
uniquely through $g_*\OO_D=\OO_V$; their analytic relations
descend through the same equality. These local factorizations
glue. Furthermore,
\[
                 k_*\OO_V=h_*\OO_D=\OO_S.
\]
Every irreducible compact curve $C\subset V$ is dominated by
a compact curve $C'\subset D$. Indeed, after base change to
the normalization of $C$, take a component dominating that curve
and a projective multisection. The projection formula gives
\[
       \delta\cdot C'
          =\deg(C'\to C)\,\epsilon\cdot C.
\]
Consequently $k$ contracts exactly the curves of
$\epsilon$-degree zero. The composite $P\to D\to V$
is a projective surjection, and $P\to S$ is projective, so
it witnesses that $k$ is Moishezon.

Conversely, suppose $k:V\to S$ is a Moishezon contraction
defined by $\epsilon$. Then $k\circ g$ is a contraction,
and the projection formula verifies its curve criterion, treating
separately curves contracted by $g$. If $Q\to V$ is a
projective surjection witnessing that $k$ is Moishezon, then
$D\times_V Q\to D$ is projective and surjective, while
$D\times_V Q\to Q\to S$ is projective. This supplies the
required witness.
\end{proof}

\begin{lemma}[Crepant transport of the multiplier closed subspace]
\label{prog:multiplier-transport}
Let $g:X\to Y$ be a projective bimeromorphic morphism of
normal compact K\"ahler spaces. Suppose effective generalized
pair data on $X$ and $Y$ have the same globally nef b-part
and are crepant over $Y$. In particular their adjoint classes
$A_X$ and $A_Y$ satisfy $A_X=g^*A_Y$ in Bott--Chern
cohomology. Write $J_X,J_Y$ for their multiplier ideals and
$N_X=V(J_X)$, $N_Y=V(J_Y)$ for their non-klt closed
subspaces. Then
\[
       g_*J_X=J_Y,\qquad R^ig_*J_X=0\quad(i>0),
       \qquad g_*\OO_{N_X}=\OO_{N_Y}.
\]
The restriction $N_X\to N_Y$ is a projective contraction.
Hence, whenever $A_Y$ is nef, the restriction of $A_X$ to
$N_X$ is endowed with a Moishezon contraction if and only if
the restriction of $A_Y$ to $N_Y$ is endowed with one.
\end{lemma}

\begin{proof}
Take a projective common log resolution $r:U\to X$ carrying
the nef datum, and put $s=g\circ r$. By crepancy there is
one real SNC divisor $C$ such that
\[
            K_U+[C]+\beta=s^*A_Y=r^*A_X,
\]
where $\beta$ is globally nef. The multiplier ideals are
\[
 J_X=r_*\OO_U(-\lfloor C\rfloor),\qquad
 J_Y=s_*\OO_U(-\lfloor C\rfloor).
\]
Their direct-image equality follows at once. For vanishing, set
$L=\OO_U(-\lfloor C\rfloor)$ and
$\Delta=C-\lfloor C\rfloor$. The pair $(U,\Delta)$
is klt, and the actual real line
\[
                  L-(K_U+\Delta)=-K_U-C
\]
has class $\beta-s^*A_Y$. It is relatively nef over $Y$
and relatively big because $s$ is bimeromorphic. The same
statements hold over $X$. Relative Kawamata--Viehweg vanishing,
in the form of \cite[Theorem~1.2]{FujinoVanishing2025} with
the first morphism the identity and the real line
$\mathcal H=-K_U-C$, therefore gives
\[
 R^is_*L=R^ir_*L=0\qquad(i>0).
\]
The real-line identity in that theorem is the displayed definition
of $\mathcal H$; its relative positivity is checked by the class
identity above. The klt condition means that there are no proper log
canonical strata to test for relative log bigness. Leray gives
$R^ig_*J_X=0$ for $i>0$. Pushing forward
$0\to J_X\to\OO_X\to\OO_{N_X}\to0$ and using
$g_*\OO_X=\OO_Y$ now proves the asserted structure-sheaf
equality. It also identifies the scheme-theoretic image with
$N_Y$, giving the projective contraction between these closed
subspaces. Apply Lemma~\ref{prog:base-emc} to conclude.
\end{proof}

\begin{remark}
The preceding lemma also transports the contraction hypothesis
through a crepant projective flip: apply it to the two contractions
to their common base, equipped with the descended adjoint and
the common nef b-data. On that base the boundary is the effective
pushforward of the source boundary, and the nef data remain on
the same carrier; the descended Bott--Chern adjoint and its
identity on a common resolution define the crepant base data
in the whole-adjoint and current-trace category used here.
All assertions concern the fixed adjoint
and its multiplier ideal. No conclusion about the multiplier
ideal of a different driving adjoint is used.
\end{remark}

\begin{corollary}[The contraction step over a lower-dimensional base]
\label{prog:base-emc-induction}
Use the hypotheses and notation of
Proposition~\ref{prog:base-presentation}. Suppose that
$A=f^*\gamma$ is nef and that its restriction to the original
non-klt closed subspace $N_o=V(J_o)$ is endowed with a
Moishezon contraction. Put
\[
                I=f_*J_1,\qquad N_*=V(I)\subset Z.
\]
Then $\gamma|_{N_*}$ is endowed with a Moishezon contraction.
If the prepared non-klt closed subspace $N_1=V(J_1)$ dominates
$Z$, then $\gamma$ and $A$ are endowed with Moishezon
contractions. Otherwise the generalized pair supplied on $Z$
has a non-klt closed subspace endowed with the contraction
defined by the restricted class of $\gamma$. In this latter
case, if $\gamma$ is nef and NQC and the assertion
$\mathsf C_{\dim Z}$ is available, then $\gamma$ and $A$
are endowed with Moishezon contractions.
\end{corollary}

\begin{proof}
The inclusion $J_o\subset J_1$ gives $N_1\subset N_o$
as closed analytic subspaces. By
Lemma~\ref{prog:base-emc}, the given contraction restricts to
$N_1$. Proposition~\ref{prog:base-presentation} identifies
\[
              f_*\OO_{N_1}=\OO_Z/I=\OO_{N_*}.
\]
Thus the restriction $N_1\to N_*$ is a projective
contraction. The restricted adjoint is the pullback of
$\gamma|_{N_*}$, so Lemma~\ref{prog:base-emc} descends
the contraction to $N_*$. Nonnegative degrees on curves of
$N_*$ follow from projective multisections and nefness of
$A$.

If $N_1$ dominates the normal irreducible space $Z$, then
$I=0$: a local function in $I$ pulls back to zero on the
dominating closed subspace and hence vanishes generically on
$Z$. Normality, and in particular reducedness, gives its
vanishing. Therefore $N_*=Z$, which proves the assertion in
this case.

If $N_1$ does not dominate $Z$, the constructed base ideal
satisfies $I\subset J_Z$. Hence
\[
                 V(J_Z)\subset V(I)=N_*
\]
as closed analytic subspaces. Restriction by
Lemma~\ref{prog:base-emc} provides the required contraction
on $V(J_Z)$, with its complete scheme structure. The
generalized pair on $Z$ has globally nef carrier data and
globally modified-big total boundary by
Proposition~\ref{prog:base-presentation}. When the stated
nef and NQC conditions hold, these facts are precisely the
hypotheses of $\mathsf C_{\dim Z}$, which gives the
contraction on $Z$. Finally Lemma~\ref{prog:base-emc}
pulls it back through the projective contraction $f$ to $X$.
The case of an empty non-klt closed subspace has the same
interpretation, with its contraction condition vacuous.
\end{proof}
